\section{The exact antiperiodic spectrum}

\begin{lemma}\label{lem:poly}
For $|z|=1$ and $s=z+z^{-1}$,
\begin{align}
 \det(xI_8-H(z))=P(x,s)
  &:=x^8-16x^6+80x^4-128x^2+38 \notag\\
  &\quad+s(-2x^4+16x^2-13)+s^2.
 \label{eq:poly}
\end{align}
Its eight roots, with all signs independently chosen, are
\begin{equation}
 \pm\sqrt{4\pm\sqrt{8+s\pm\sqrt{26-3s}}}.
 \label{eq:all-roots}
\end{equation}
They are all real, nonzero, and distinct for every $-2\le s\le2$.
\end{lemma}
\begin{proof}
Expanding the determinant of \eqref{eq:fiber} as a Laurent polynomial
in $z$ gives
\[
 x^8-16x^6+80x^4-128x^2+40
 +(z+z^{-1})(-2x^4+16x^2-13)+(z^2+z^{-2}).
\]
Using $z^2+z^{-2}=s^2-2$ proves \eqref{eq:poly}.
This is a finite polynomial identity and can also be checked by the exact
symbolic determinant in the accompanying verifier.

Put $y=x^2$. A translation gives
\begin{equation}
 P(x,s)=(y-4)^4-(16+2s)(y-4)^2+s^2+19s+38.
 \label{eq:centered}
\end{equation}
Solving this as a quadratic in $(y-4)^2$ gives
\[
 u_\pm(s)=8+s\pm\sqrt{26-3s}.
\]
The function $s^2+19s+38$ is increasing on $[-2,2]$ and has value $4$
at $-2$. Since $8+s>0$, the identity
\[
 (8+s)^2-(26-3s)=s^2+19s+38\ge4
\]
proves $u_->0$. Also $u_+>u_-$, and
\[
 u_+'(s)=1-\frac3{2\sqrt{26-3s}}>0,
 \qquad u_+(s)\le10+2\sqrt5<16.
\]
Consequently
\[
 0<4-\sqrt{u_+}<4-\sqrt{u_-}
   <4+\sqrt{u_-}<4+\sqrt{u_+}.
\]
These are the four distinct positive roots in $y$. Taking their positive
and negative square roots gives exactly \eqref{eq:all-roots}.
\end{proof}

\begin{proof}[Proof of Theorem~\ref{thm:main}]
Lemma~\ref{lem:poly} implies
\[
 \rho(H(z))=R(s),\qquad s=z+z^{-1}.
\]
The derivative computation in that lemma and strict monotonicity of the
outer square roots show that $R$ increases strictly on $[-2,2]$.
For positive holonomy, the grid $z^m=1$ contains $z=1$, so its maximal
value of $s$ is $2$. For negative holonomy, the grid is
\[
 z_j=\exp\!\left(\frac{(2j+1)\pi\mathrm i}{m}\right),
 \qquad 0\le j<m,
\]
and its maximal value of $s$ is $2\cos(\pi/m)$. Lemma~\ref{lem:bloch}
therefore proves both equalities in \eqref{eq:main}.
For every finite $m\ge1$, $\cos(\pi/m)<1$, which proves the strict
inequality. Lemma~\ref{lem:gauge} shows that these are exactly the two
switching classes for the prescribed triangle signs, so the strict
comparison proves the stated constrained uniqueness. Since $A_m^-$ is
an admissible real signing, the unrestricted upper bound follows.
\end{proof}

To identify \eqref{eq:rstar} with the algebraic number used in the revised
conjecture, put
\[
 f(x)=x^4-2x^3-6x^2+12x-4.
\]
An exact expansion gives $P(x,2)=f(x)f(-x)$. On $x\ge\sqrt6$,
\[
 f''(x)=12(x^2-x-1)>0,\qquad
 f'(\sqrt6)=12\sqrt6-24>0,\qquad f(\sqrt6)=-4.
\]
Hence $f$ has exactly one root on that ray. For $x>\sqrt6$,
$f(-x)-f(x)=4x(x^2-6)>0$. The largest positive root of $P(x,2)$ is
$R(2)>\sqrt6$. It must be a root of $f$, because $f(-R(2))=0$ would
force $f(R(2))<0$ and then a still larger positive root of $f$, contrary
to maximality. Thus $R(2)$ is the largest real root of $f$.

\begin{corollary}\label{cor:asymptotic}
The negative-holonomy radii increase strictly to $\rstar$ as $m$ increases,
and
\begin{equation}
 \rstar-\rho(A_m^-)
 =\frac{\pi^2}{4\rstar\sqrt{10+2\sqrt5}}
   \left(1-\frac3{4\sqrt5}\right)m^{-2}+O(m^{-4}).
 \label{eq:asymptotic}
\end{equation}
\end{corollary}
\begin{proof}
The sequence $2\cos(\pi/m)$ increases strictly to $2$, and $R$ is strictly
increasing and smooth near $2$. Since
\[
 2\cos(\pi/m)=2-\pi^2m^{-2}+O(m^{-4}),
 \qquad
 R'(2)=\frac{1-3/(4\sqrt5)}{4\rstar\sqrt{10+2\sqrt5}},
\]
Taylor's theorem proves \eqref{eq:asymptotic}.
\end{proof}

There is no parity exception: when $m$ is odd, the allowed phase $z=-1$
has $s=-2$ and cannot recover the omitted maximum at $s=2$.
At $m=1$ the single fiber is $H(-1)$, with
$\rho(A_1^-)=\sqrt{6+\sqrt2}$.
